\clearpage
\SourceLeaf{177}\begin{center}
{\Large\textls{ANZEIGE.}}
\end{center}

\bigskip
\hrule
\begin{center}
Göttingische gelehrte Anzeigen. \quad 1816 April 20.
\end{center}
\hrule

\bigskip
\begin{center}
{\large\textls{Stuttgart.}}
\end{center}

Typis J. F. Steinkopf: \emph{Commentatio in primum elementorum Euclidis librum,
qua veritatem geometriae principiis ontologicis niti evincitur, omnesque
  propositiones, axiomatum geometricorum loco habitae, demonstrantur. Auctore}
  J. C. \textsc{Schwab}, \emph{Regi Württembergiae a consiliis aulicis
secretioribus, academiae scientiarum Petropolitanae, Berolinensis et Harlemensis
Sodali.} 1814. 65 Seiten in Octav.

\begin{center}
{\large\textls{Mainz.}}
\end{center}

Auf Kosten des Verfassers und in Commission bei Florian Kupferberg:
\emph{Vollständige Theorie der Parallellinien. Nebst einem Anhange, in welchem
der erste Grundsatz zur Technik der geraden Linie angegeben wird. Herausgegeben
von} \textsc{Matthias Metternich}, \emph{Doctor der Philosophie, Professor der
Mathematik, Mitglied der gelehrten Gesellschaft nützlicher Wissenschaften zu
Erfurt.} 1815. 44 Seiten in Octav.

Es wird wenige Gegenstände im Gebiete der Mathematik geben, über welche so viel
geschrieben wäre, wie über die Lücke im Anfange der Geometrie bei Begründung der
Theorie der Parallellinien. Selten vergeht ein Jahr, wo nicht irgend ein neuer
Versuch zum Vorschein käme, diese Lücke auszufüllen, ohne dass wir doch, wenn wir
ehrlich und offen reden wollen, sagen könnten, dass wir im Wesentlichen irgend
weiter gekommen wären, als Euklides vor 2000 Jahren war. Ein solches
aufrichtiges und unumwundenes Ge%
